% !TeX root = ../main.tex

\setproblem{I}{\ProblemTitleI}
\section{\texorpdfstring{\pletter. \ptitle}{문제 \pletter}}

\begin{frame}[label=problem-I]
    \sectiontitle{\pletter}{\ptitle}
    \sectionmeta{
        출제자: \idtext{areian13}\\
        태그: \tagtext{\#dp, \#combinatorics}\\
        예상 난이도: \tiertext{Silver3}
    }
    \begin{itemize}
        \item 제출 12번, 정답 7명 (정답률 58.3\%)
        \item 처음 푼 사람: ㅈㄱㅇ (00:24:39)
    \end{itemize}
\end{frame}

\begin{frame}{\textbf{\pletter}. \ptitle}
\begin{itemize}
    \item 같은 종류의 진공관이 $k$개 연속해 있다고 해 봅시다.
    \item 마지막 진공관을 따로 포장하면 앞의 $k-1$개를, 짝지으면 앞의 $k-2$개를 포장하면 됩니다.
    \item 즉, 피보나치 점화식 $fib[k]=fib[k-1]+fib[k-2]$입니다.
    \item $90$까지 미리 계산해 두고, 연속한 진공관 묶음마다 해당하는 값을 곱하면 됩니다.
\end{itemize}
\end{frame}
